% Part: proof-theory
% Chapter: cut-elimination
% Section: midsequent

\documentclass[../../../include/open-logic-section]{subfiles}

\begin{document}

\olfileid{pt}{cut}{mid}

\olsection{د منځني سېکوېنټ قضيه او د هېربراند قضيه}

د پرېکون لرې کولو د قضيې يوه مهمه پايله د منځني سېکوېنټ
قضيه ده۔ دا وايي: که يو !!a{proof}~$\pi$ د~$\Gamma \Sequent \Delta$
وي چې پکښې هر !!{formula} په مقدم کميتي بڼه وي، نو يو !!a{proof}~$\pi'$
شته چې يو سېکوېنټ $S \ident \Pi \Sequent
\Lambda$ (\emph{منځنے سېکوېنټ}) لري، داسې چې په~$\pi'$ کښې د~$S$
پاس ټول استنتاجونه گزاره‌يي او د~$S$ لاندې ټول استنتاجونه د کميت
اچوونکو وي۔ (!!^a{formula}~$!A$ ته \emph{په مقدم کميتي بڼه}
وايي که دا بڼه ولري:
\[\mathsf{Q}_1x_1\dots\mathsf{Q}_nx_n\,!B(x_1,\dots,x_n)\]
چې هر $\mathsf{Q}_i$ پکښې يا $\lforall$ يا~$\lexists$ وي او وروستنے
متن له کميت اچوونکو خالي وي۔ \textbf{سرچينه‌يي سمون OLCMP-126:}
د مقدم کميتي بڼې په تعريف کښې د وروستي متن د کميت نشتوالي شرط
څرګند شوے دے۔) د منځني سېکوېنټ د قضيې يوه مهمه پايله د هېربراند
قضيه ده۔ عام حالت ئې بيانول لږ سخت دي، خو په پراخه توګه دا بڼه
لري: د لومړۍ درجې منطق د هرې !!{provable} !!{sentence}~$!A$ دپاره
يوه گزاره‌يي توتولوژي~$!A'$ شته چې ترې $!A$ !!{prove}d کېدے شي۔
دا ئې د هغو !!{sentence}s دپاره بڼه ده چې د~$\lexists[x][!B(x)]$
په شکل وي، چې $!B$ له کميت اچوونکو خالي وي:


\begin{thm}[د هېربراند قضيه]
فرض کړئ چې $!B(x)$ له کميت اچوونکو خالي دے او $\Proves 
\lexists[x][!B(x)]$۔ نو ترمونه~$t_1$، \dots، $t_n$ شته داسې چې
$\Proves !B(t_1) \lor \dots \lor !B(t_n)$۔
\end{thm}

!!{formula}~$!B(t_1) \lor \dots \lor !B(t_n)$ ته د~$\lexists[x][!B(x)]$
\emph{د هېربراند فصل} وايي۔ څرنګه چې $!B$ له کميت اچوونکو خالي
دے، دا فصل هم له هغو خالي دے۔ نو که دا !!{provable} وي، په
بې‌پرېکونه ثبوت !!{provable} دے، او په نتيجه کښې په داسې
!!a{proof} چې يوازې گزاره‌يي استنتاجونه کاروي۔ ځکه نو توتولوژي ده۔

د هېربراند د قضيې سخته برخه دا ښودل دي چې د هر !!{provable}
!!{formula}~$\lexists[x][!B(x)]$ دپاره د هېربراند يو فصل شته۔
معکوس ئې، يعنې که $\lexists[x][!B(x)]$ د هېربراند فصل ولري نو
ثبوت لري، اسان دے۔ په حقيقت کښې د هېربراند فصل، $\lexists[x][!B(x)]$
لازموي۔ مثلاً په~$\Log{G3c}$ کښې د $n=2$ حالت دا !!a{proof} دے:
\[
\Axiom$!B(t_1) \fCenter \lexists[x][!B(x)], !B(t_1), !B(t_2)$
\Axiom$!B(t_2) \fCenter \lexists[x][!B(x)], !B(t_1), !B(t_2)$
\RightLabel{\LeftR{\lor}}
\BinaryInf$!B(t_1) \lor !B(t_2) \fCenter \lexists[x][!B(x)], !B(t_1), !B(t_2)$
\RightLabel{\RightR{\lexists}}
\UnaryInf$!B(t_1) \lor !B(t_2) \fCenter \lexists[x][!B(x)], !B(t_2)$
\RightLabel{\RightR{\lexists}}
\UnaryInf$!B(t_1) \lor !B(t_2) \fCenter \lexists[x][!B(x)]$
\DisplayProof
\]

د $\Sequent
\lexists[x][!B(x)]$ له !!a{proof}~$\pi$ څخه د هېربراند د فصل
راايستلو دپاره هغه حالت وګورئ چې په يوه بې‌پرېکونه !!{proof}~$\pi$
کښې ټول $\RightR{\lexists}$ استنتاجونه په پای کښې وي:
\[
\AxiomC{}
\RightLabel{$\pi_1$}
\Deduce$ \fCenter \lexists[x][!B(x)], !B(t_1), \dots, !B(t_n)$
\RightLabel{\RightR{\lexists}}
\UnaryInf$ \fCenter \lexists[x][!B(x)], !B(t_2), \dots, !B(t_n)$
\RightLabel{$(n-1) \times \RightR{\lexists}$}
\Deduce$ \fCenter \lexists[x][!B(x)]$
\DisplayProof
\]
که په~$\pi$ کښې دا ټول هغه $\RightR{\lexists}$ استنتاجونه وي چې
$\lexists[x][!B(x)]$ پکښې فعال وي، نو په~$\pi_1$ کښې نور کميتي
استنتاجونه نشته۔ \textbf{سرچينه‌يي سمون OLCMP-127:} دلته د
هماغه رسم موجود فرعي ثبوت ته رجوع سمه شوې ده۔ ځکه چې $!B(x)$
نور کميت اچوونکي نۀ لري او پای سېکوېنټ په کيڼ اړخ کښې هېڅ
!!{formula}s نۀ لري۔ په نورو ټکو، $ \fCenter \lexists[x][!B(x)], !B(t_1),
\dots, !B(t_n)$ د~$\pi$ منځنے سېکوېنټ دے۔ همدارنګه، د منځني
سېکوېنټ پاس کميتي استنتاجونه نشته، نو پاس ټول سېکوېنټونه
$\lexists[x][!B(x)]$ لري او دا نۀ فعال دي او نۀ اصلي۔ ځکه دا له
!!{proof}~$\pi_1$ څخه لرې کولے شو او د~$\fCenter !B(t_1), \dots, !B(t_n)$
يو !!a{proof} تر لاسه کوو؛ بيا د $\RightR{\lor}$ د $n-1$
استعمالونو په وسيله د هېربراند د فصل~$\Sequent !B(t_1) \lor \dots \lor !B(t_n)$
يو !!a{proof} تر لاسه کوو۔

ستونزه دا ده چې فرض کولے نۀ شو چې د $\Sequent \lexists[x][!B(x)]$
يو بې‌پرېکونه !!{proof} هم خپل ټول \RightR{\lexists} استنتاجونه
په پای کښې په يوه ډله کښې لري۔ د دې \RightR{\lexists} استنتاجونو
تر منځ داسې استنتاجونه کېدے شي چې ځينې $!B(t_i)$ پکښې اصلي وي۔
ځکه نو د منځني سېکوېنټ قضيې ته اړتيا ده۔

\begin{thm}[د منځني سېکوېنټ قضيه]\ollabel{thm:midsequent}
که د $\Gamma \Sequent \Delta$ يو \Log{G3c}-!!{proof}~$\pi$ وي، چې
په $\Gamma$ او $\Delta$ کښې ټول !!{formula}s په مقدم کميتي بڼه
وي، نو يو !!a{proof}~$\pi'$ شته چې سېکوېنټ $\Pi \Sequent \Lambda$
لري، داسې چې لاندې ترې ټول استنتاجونه کميتي او پاس ترې ټول
استنتاجونه گزاره‌يي وي۔
\end{thm}

\begin{proof}
د پرېکون لرې کولو په وسيله لومړے بې‌پرېکونه ثبوت اخلو۔ په~$\pi$
کښې د يوه کميتي استنتاج \emph{ترتيبي شمېر} د لاندې گزاره‌يي
استنتاجونو شمېر تعريفوو، او د $\pi$ ترتيبي شمېر~$o(\pi)$ د ټولو
کميتي استنتاجونو د ترتيبي شمېرونو مجموعه ده۔ که $o(\pi) = 0$
وي، نو د هېڅ کميتي استنتاج لاندې گزاره‌يي استنتاج نشته او کار
ختم دے: که کميتي استنتاج نۀ وي، پای سېکوېنټ منځنے سېکوېنټ
اخلو؛ که وي، د ټولو کميتي استنتاجونو يوه مقدمه ده، نو يو واحد
تر ټولو پاسنے کميتي استنتاج شته۔ د هغه مقدمه منځنے سېکوېنټ ده۔
\textbf{سرچينه‌يي سمون OLCMP-129:} د فرعي فارمول د خاصيت د
کارولو دپاره بې‌پرېکونه ثبوت اخستل او د کميتي استنتاج د نشتوالي
بنسټيز حالت څرګند شوي دي۔

که $o(\pi) > 0$ وي، د $>0$ ترتيبي شمېر لرونکو کميتي
استنتاجونو تر ټولو لاندې يو غوره کړئ۔ څرنګه چې ترتيبي شمېر
ئې $>0$ دے، لاندې ترې يو گزاره‌يي استنتاج ضرور شته۔ دا چې تر
ټولو لاندې يو دے، پايله ئې د بل کميتي استنتاج مقدمه کېدے نۀ
شي: ګنې د هغه دويم لاندې کميتي استنتاج لاندې به هم يو گزاره‌يي
استنتاج وے۔ د کميتي استنتاج او د لاندې گزاره‌يي استنتاج له
ډول سره سم (ډېر!) حالتونه بېلوو۔ څرنګه چې پای سېکوېنټ يوازې
د مقدم کميتي بڼې فارمولونه لري، د کميتي استنتاج اصلي فارمول
په ورپسې گزاره‌يي استنتاج کښې فعال کېدے نۀ شي۔ ګنې د هغه
استنتاج د اصلي فارمول په يوه گزاره‌يي رابط کښې به کميت اچوونکے
راشي۔ د فرعي !!{formula} د خاصيت له مخې به دا د پای سېکوېنټ
د يوې !!a{formula} فرعي !!{formula} وے۔ نو د کميتي استنتاج
اصلي !!{formula} د لاندې گزاره‌يي استنتاج د سياق يو !!a{element} دے۔

دا دوه بېلګې دي:

لومړے فرض کړئ چې يو $\RightR{\lforall}$ د $\LeftR{\lif}$ په
ښۍ مقدمه ختمېږي۔ نو د~$\pi$ اړوند فرعي !!{proof} په فرعي
!!{proof}~$\pi_1$ ختمېږي:
\[
\AxiomC{}
\RightLabel{$\pi_2$}
\Deduce$\Gamma' \fCenter \Delta', \lforall[x][!A(x)], !B$
\AxiomC{}
\RightLabel{$\pi_3$}
\Deduce$!C, \Gamma' \fCenter \Delta', !A(c)$
\RightLabel{\RightR{\lforall}}
\UnaryInf$!C, \Gamma' \fCenter \Delta', \lforall[x][!A(x)]$
\RightLabel{\LeftR{\lif}}
\BinaryInf$!B \lif !C, \Gamma' \fCenter \Delta', \lforall[x][!A(x)]$
\DisplayProof
\]
فرض کولے شو چې $\pi$ منظم دے، نو ځانګړے متغير~$c$ له~$\pi_3$
بهر نۀ راځي؛ په ځانګړې توګه په $!B
\lif !C$ او په~$\pi_2$ کښې نۀ راځي۔ اوس ثبوت~$\pi_1'$ وګورئ:
\[
\AxiomC{}
\RightLabel{$\pi_2'$}
\Deduce$\Gamma' \fCenter \Delta', \lforall[x][!A(x)], !A(c), !B$
\AxiomC{}
\RightLabel{$\pi_3$}
\Deduce$!C, \Gamma' \fCenter \Delta', \lforall[x][!A(x)], !A(c)$
\RightLabel{\LeftR{\lif}}
\BinaryInf$!B \lif !C, \Gamma' \fCenter \Delta', \lforall[x][!A(x)], !A(c)$
\RightLabel{\RightR{\lforall}}
\UnaryInf$!B \lif !C, \Gamma' \fCenter \Delta', \lforall[x][!A(x)], \lforall[x][!A(x)]$
\DisplayProof
\]
چې $\pi_2'$ له~$\pi_2$ څخه په ښي اړخ کښې د~$!A(c)$ په
کمزوري کولو تر لاسه کېږي۔ (دا هېڅ استنتاج نۀ زياتوي، په
ځانګړې توګه داسې کميتي استنتاج نۀ زياتوي چې ترتيبي شمېر
بدلولے شي۔ په~$\pi_2$ کښې د ځانګړے متغير شرطونه هم نۀ ماتوي،
ځکه په~$!A(c)$ کښې هېڅ ثابت په~$\pi_2$ کښې د ځانګړے متغير په
توګه نۀ کارېږي۔) پر \RightR{\lforall} د ځانګړے متغير شرط لا هم
پوره دے، ځکه $c$ په~$!B \lif !C$ کښې نۀ راځي۔

د انقباض د منلو وړتيا کاروو چې د $!B \lif !C, \Gamma' \fCenter \Delta',
\lforall[x][!A(x)]$ يو !!a{proof}~$\pi_1''$ تر لاسه کړو۔ د
$\lforall[x][!A(x)]$ دپاره د $\RightR{\Contraction}$ د منلو
وړتيا ثبوت او په هغه کښې د $\RightR{\lforall}$ د معکوسولو وړتيا
ثبوت د هېڅ استنتاج ترتيبي شمېر نۀ بدلولو: تر ډېره ئې استنتاجونه
لرې کړي وو۔ نو په~$\pi_1''$ کښې کميتي استنتاجونه تر~$\pi_1$
زيات نۀ دي، او په~$\pi_1''$ کښې د هر کميتي استنتاج لاندې د
گزاره‌يي استنتاجونو شمېر په~$\pi_1$ کښې د ورته استنتاج د
لاندې استنتاجونو هومره دے؛ پرته له وروستي \RightR{\lforall}:
په~$\pi$ کښې يو \LeftR{\lif} لاندې ترې و، چې په~$\pi_1''$ کښې
مو پاس ترې يووړ۔ په $\pi$ کښې $\pi_1$ په $\pi_1''$ بدل کړئ۔
د نوي !!{proof} ترتيبي شمېر لږ دے، ځکه لږ تر لږه د
$\RightR{\lforall}$ استنتاج ترتيبي شمېر (لږ تر لږه)~$1$ کم
شوے دے۔ اوس استقرايي فرضيه وکاروئ۔

هغه حالتونه چې کميتي استنتاج \RightR{\lexists} وي نور هم
اسان دي، ځکه انقباض نۀ غواړي او د ځانګړے متغير شرط ته اندېښنه
نۀ وي۔ مثلاً فرض کړئ چې له $\RightR{\lexists}$ وروسته
$\RightR{\land}$ راځي (دا ځل د کيڼې مقدمې په توګه)۔ اړوند فرعي !!{proof}
$\pi_1$ دے:
\[
\AxiomC{}
\RightLabel{$\pi_2$}
\Deduce$\Gamma' \fCenter \Delta', \lexists[x][!A(x)], !A(t), !B$
\RightLabel{\RightR{\lexists}}
\UnaryInf$\Gamma' \fCenter \Delta', \lexists[x][!A(x)], !B$
\AxiomC{}
\RightLabel{$\pi_3$}
\Deduce$\Gamma' \fCenter \Delta', \lexists[x][!A(x)], !C$
\RightLabel{\RightR{\land}}
\BinaryInf$\Gamma' \fCenter \Delta', \lexists[x][!A(x)], !B \land !C$
\DisplayProof
\]
\textbf{سرچينه‌يي سمون OLCMP-128:} د رسم له مقدم سياقه
اضافي د فارمول نښه لرې شوه؛ اصلي سياق عين پاتې دے۔
په~$\pi$ کښې $\pi_1$ په~$\pi_1'$ بدل کړئ:
\[
\AxiomC{}
\RightLabel{$\pi_2$}
\Deduce$\Gamma' \fCenter \Delta', \lexists[x][!A(x)], !A(t), !B$
\AxiomC{}
\RightLabel{$\pi_3'$}
\Deduce$\Gamma' \fCenter \Delta', \lexists[x][!A(x)], !A(t), !C$
\RightLabel{\RightR{\land}}
\BinaryInf$\Gamma' \fCenter \Delta', \lexists[x][!A(x)], !A(t), !B \land !C$
\RightLabel{\RightR{\lexists}}
\UnaryInf$\Gamma' \fCenter \Delta', \lexists[x][!A(x)], !B \land !C$
\DisplayProof
\]
چې $\pi_3'$ له $\pi_3$ څخه په ښي اړخ کښې د~$!A(t)$ په
کمزوري کولو تر لاسه کېږي۔ دا کمزوري کول يوازې د~$\pi_3$ د
هر سېکوېنټ ښي اړخ ته $!A(t)$ زياتوي، نو د هېڅ کميتي استنتاج
ترتيبي شمېر نۀ بدلوي۔ په~$\pi'$ کښې د کميتي استنتاجونو ترتيبي
شمېرونه د~$\pi$ په شان دي، پرته له هغه \RightR{\lexists}
استنتاج څخه چې له~\RightR{\land} سره بدل شوے او ترتيبي شمېر
ئې~$1$ کم شوے دے۔ استقرايي فرضيه کارولے شو۔
\end{proof}

\begin{prob}
د \olref{thm:midsequent} په ثبوت کښې هغه حالتونه بيان کړئ چې
\LeftR{\lexists} د~\RightR{\lif} په مقدمه او \LeftR{\lforall}
د~\LeftR{\lor} په کيڼه مقدمه ختمېږي۔
\end{prob}

\begin{proof}[د هېربراند د قضيې ثبوت]
فرض کړئ چې $\pi$ په $\Sequent \lexists[x][!A(x)]$ ختمېږي۔ د
\olref{thm:midsequent} له مخې $\pi$ په داسې !!a{proof}~$\pi'$
بدلولے شو چې منځنے سېکوېنټ~$S$ ولري، او هغه ضرور دا بڼه لري:
\[\Sequent
\lexists[x][!A(x)], !A(t_1), \dots, !A(t_n).\]
څرنګه چې د $S$ پاس ټول استنتاجونه گزاره‌يي دي، $\lexists[x][!A(x)]$
د~$S$ پاس د کوم استنتاج اصلي فارمول کېدے نۀ شي۔ دا چې اټومي
نۀ دے، په بديهيه کښې هم اصلي کېدے نۀ شي۔ څرنګه چې
$\lexists[x][!A(x)]$ د $!A(t_1)$ اصلي فرعي فارمول کېدے نۀ شي
(ياد مو وي چې $!A(x)$ له کميت اچوونکو خالي دے)، د~$S$ پاس
فعال هم کېدے نۀ شي۔ ځکه نو په~$S$ ختمېدونکي فرعي !!{proof}
کښې د $\lexists[x][!A(x)]$ لرې کول د $\Sequent !A(t_1), \dots, !A(t_n)$
يو کره !!{proof} ورکوي، چې ترې د $\RightR{\lor}$ د $n-1$
استنتاجونو په وسيله د $\Sequent !A(t_1) \lor \dots \lor !A(t_n)$
يو !!a{proof} تر لاسه کېدے شي۔
\end{proof}

% \begin{prob}
% Find a Herbrand disjunction of $\lexists[x][\lforall[y][(!A(x) \lif
% !A(y))]]$ by finding a cut-free !!{proof} of $\Sequent
% \lexists[x][\lforall[y][(!A(x) \lif !A(y))]]$ and applying the
% transformation in \olref{thm:midsequent} to find a midsequent (if
% necessary).
% \end{prob}

\end{document}
